\documentclass[aps,preprint]{revtex4}%
\usepackage{amsfonts}
\usepackage{amsmath}
\usepackage{amssymb}
\usepackage{graphicx}%
\setcounter{MaxMatrixCols}{30}
%TCIDATA{OutputFilter=latex2.dll}
%TCIDATA{Version=5.50.0.2953}
%TCIDATA{CSTFile=revtex4.cst}
%TCIDATA{Created=Monday, November 19, 2007 08:47:41}
%TCIDATA{LastRevised=Thursday, January 10, 2008 15:25:16}
%TCIDATA{<META NAME="GraphicsSave" CONTENT="32">}
%TCIDATA{<META NAME="SaveForMode" CONTENT="1">}
%TCIDATA{BibliographyScheme=Manual}
%TCIDATA{<META NAME="DocumentShell" CONTENT="Articles\SW\REVTeX 4">}
%BeginMSIPreambleData
\providecommand{\U}[1]{\protect\rule{.1in}{.1in}}
%EndMSIPreambleData
\newtheorem{theorem}{Theorem}
\newtheorem{acknowledgement}[theorem]{Acknowledgement}
\newtheorem{algorithm}[theorem]{Algorithm}
\newtheorem{axiom}[theorem]{Axiom}
\newtheorem{claim}[theorem]{Claim}
\newtheorem{conclusion}[theorem]{Conclusion}
\newtheorem{condition}[theorem]{Condition}
\newtheorem{conjecture}[theorem]{Conjecture}
\newtheorem{corollary}[theorem]{Corollary}
\newtheorem{criterion}[theorem]{Criterion}
\newtheorem{definition}[theorem]{Definition}
\newtheorem{example}[theorem]{Example}
\newtheorem{exercise}[theorem]{Exercise}
\newtheorem{lemma}[theorem]{Lemma}
\newtheorem{notation}[theorem]{Notation}
\newtheorem{problem}[theorem]{Problem}
\newtheorem{proposition}[theorem]{Proposition}
\newtheorem{remark}[theorem]{Remark}
\newtheorem{solution}[theorem]{Solution}
\newtheorem{summary}[theorem]{Summary}
\newenvironment{proof}[1][Proof]{\noindent\textbf{#1.} }{\ \rule{0.5em}{0.5em}}
\begin{document}
\preprint{HEP/123-qed}
\title[Short title for running header]{Long title that appears on the title page}
\author{First Author}
\affiliation{My Institution}
\author{Second Author}
\affiliation{My Institution}
\author{Third Author}
\affiliation{Other Institution}
\keywords{one two three}
\pacs{PACS number}

\begin{abstract}
Shell document for REV\TeX{} 4.

\end{abstract}
\volumeyear{year}
\volumenumber{number}
\issuenumber{number}
\eid{identifier}
\date[Date text]{date}
\received[Received text]{date}

\revised[Revised text]{date}

\accepted[Accepted text]{date}

\published[Published text]{date}

\startpage{101}
\endpage{102}
\maketitle
\tableofcontents


\section{Geminal Power States}

The general expressions for the AGP SORDO and FORDO\ are
\begin{align*}
D^{2}\left(  \left\vert g^{\frac{N}{2}}\right\rangle \left\langle g^{\frac
{N}{2}}\right\vert \right)   &  =\left(  S_{\frac{N}{2}}\right)  ^{-1}\left\{
\sum_{1\leq j\leq s}n_{j}S_{\frac{N}{2}-1}\left(  \hat{\jmath}\right)
\left\vert \psi_{j}\psi_{j+s}\right\rangle \left\langle \psi_{j}\psi
_{j+s}\right\vert \right. \\
&  +\sum_{1\leq j\neq k\leq s}c_{j}c_{k}e^{i\left(  \theta_{j}-\theta
_{k}\right)  }S_{\frac{N}{2}-1}\left(  \hat{\jmath}\hat{k}\right)  \left\vert
\psi_{j}\psi_{j+s}\right\rangle \left\langle \psi_{k}\psi_{k+s}\right\vert \\
&  \left.  +\sum_{\substack{1\leq j<k\leq r\\k\neq j+s}}n_{j}n_{k}S_{\frac
{N}{2}-2}\left(  \hat{\jmath}\hat{k}\right)  \left\vert \psi_{j}\psi
_{k}\right\rangle \left\langle \psi_{j}\psi_{k}\right\vert \right\}  .
\end{align*}
and
\[
D^{1}\left(  \left\vert g^{\frac{N}{2}}\right\rangle \left\langle g^{\frac
{N}{2}}\right\vert \right)  =\left(  S_{\frac{N}{2}}\right)  ^{-1}\sum_{1\leq
j\leq s}n_{j}S_{\frac{N}{2}-1}\left(  \hat{\jmath}\right)  \left\{  \left\vert
\psi_{j}\right\rangle \left\langle \psi_{j}\right\vert +\left\vert \psi
_{j+s}\right\rangle \left\langle \psi_{j+s}\right\vert \right\}
\]
repectively, where $n_{j}=c_{j}^{2}$. The energy of an AGP\ state is given by
\begin{align*}
&  \mathcal{E}\left(  \mathbf{n},\mathbf{\psi,\theta}\right)  =\frac
{\mathcal{H}\left(  \mathbf{n},\mathbf{\psi}\right)  }{S_{\frac{N}{2}}\left(
\mathbf{n}\right)  }=\\
&  \left\{  \sum_{1\leq j\leq s}n_{j}S_{\frac{N}{2}-1}\left(  \hat{\jmath
}\right)  \left\{  \overset{}{\left\langle \psi_{j}\left\vert h_{1}\psi
_{j}\right.  \right\rangle +\left\langle \psi_{j+s}\left\vert h_{1}\psi
_{j+s}\right.  \right\rangle }+\,\left\langle \psi_{j}\psi_{j+s}\right\vert
\left\vert \psi_{j}\psi_{j+s}\right\rangle \right\}  \right. \\
&  +\sum_{1\leq j\neq k\leq s}\left\{  c_{j}c_{k}e^{i\left(  \theta_{j}%
-\theta_{k}\right)  }S_{\frac{N}{2}-1}\left(  \hat{\jmath}\hat{k}\right)
\left\langle \psi_{j}\psi_{j+s}\right\vert \left\vert \psi_{k}\psi
_{k+s}\right\rangle \right. \\
&  \qquad\qquad\qquad\left.  \left.  \left.  +\frac{1}{2}\overset{}{n_{j}%
n_{k}S_{\frac{N}{2}-2}\left(  \hat{\jmath}\hat{k}\right)  \left\langle
\psi_{j}\psi_{k}|||\psi_{j}\psi_{k}\right\rangle }\right\}  \right\}  \right/
S_{\frac{N}{2}}\left(  \mathbf{n}\right)  ,
\end{align*}
where
\begin{align*}
\left\langle \psi_{j}\psi_{k}|||\psi_{j}\psi_{k}\right\rangle  &
=\left\langle \psi_{j}\psi_{k}||\psi_{j}\psi_{k}\right\rangle +\left\langle
\psi_{j}\psi_{k+s}||\psi_{j}\psi_{k+s}\right\rangle \\
&  +\left\langle \psi_{j+s}\psi_{k}||\psi_{j+s}\psi_{k}\right\rangle
+\left\langle \psi_{j+s}\psi_{k+s}||\psi_{j+s}\psi_{k+s}\right\rangle
\end{align*}
and
\begin{align*}
\left\langle \psi_{j}\psi_{k}||\psi_{l}\psi_{m}\right\rangle  &  =\int\left\{
\psi_{j}^{\ast}\left(  \mathbf{r}_{1}\mathbf{,\xi}_{1}\right)  \psi_{k}^{\ast
}\left(  \mathbf{r}_{2}\mathbf{,\xi}_{2}\right)  \frac{1}{\left\Vert
\mathbf{r}_{1}-\mathbf{r}_{2}\right\Vert }\psi_{l}\left(  \mathbf{r}%
_{1}\mathbf{,\xi}_{1}\right)  \psi_{m}\left(  \mathbf{r}_{2}\mathbf{,\xi}%
_{2}\right)  \right. \\
&  -\left.  \psi_{j}^{\ast}\left(  \mathbf{r}_{1}\mathbf{,\xi}_{1}\right)
\psi_{k}^{\ast}\left(  \mathbf{r}_{2}\mathbf{,\xi}_{2}\right)  \frac
{1}{\left\Vert \mathbf{r}_{1}-\mathbf{r}_{2}\right\Vert }\psi_{l}\left(
\mathbf{r}_{2}\mathbf{,\xi}_{2}\right)  \psi_{m}\left(  \mathbf{r}%
_{1}\mathbf{,\xi}_{1}\right)  \right\}  d\mathbf{r}_{1}d\mathbf{r}_{2}d\xi
_{1}d\xi_{2}.
\end{align*}
The symmetric functions are given by
\begin{align*}
S_{m}  &  \equiv S_{m}\left(  \mathbf{n}\right)  =\sum_{1\leq j_{1}%
<\cdots<j_{m}\leq s}n_{j_{1}}\cdots n_{j_{m}}\\
S_{m}  &  =0,\,m<0\\
S_{0}  &  =1\\
S_{1}  &  =\sum_{1\leq j_{1}\leq s}n_{j_{1}}\\
S_{1}\left(  \hat{k}\right)   &  =\sum_{1\leq j_{1}\neq k\leq s}n_{j_{1}}\\
S_{1}\left(  \hat{k}\hat{l}\right)   &  =\sum_{1\leq j_{1}\neq k,l\leq
s}n_{j_{1}}\\
S_{2}  &  =\sum_{1\leq j_{1}<j_{2}\leq s}n_{j_{1}}n_{j_{2}}\\
S_{2}\left(  \hat{k}\right)   &  =\sum_{\substack{1\leq j_{1}<j_{2}\leq
s\\j_{1},j_{2}\neq k}}n_{j_{1}}n_{j_{2}}.
\end{align*}


\section{Example States}

We will consider low spin geminals
\[
S_{0}\left\vert g\right\rangle =0
\]
formed from spin orbitals
\[
S_{0}\left\vert \psi\right\rangle =\pm\frac{1}{2}\left\vert \psi\right\rangle
\]
that produce the same FORDO, starting with the geminal%

\[
\left\vert g_{+}\right\rangle =\left\vert \psi_{1}\psi_{3}\right\rangle
+\left\vert \psi_{2}\psi_{4}\right\rangle =\left\vert \varphi_{1}\alpha
\varphi_{3}\beta\right\rangle +\left\vert \varphi_{2}\alpha\varphi_{4}%
\beta\right\rangle
\]%
\begin{align}
&  =\frac{1}{\sqrt{2}}\left\{  \varphi_{1}\left(  1\right)  \alpha\left(
1\right)  \otimes\varphi_{3}\left(  2\right)  \beta\left(  2\right)
-\varphi_{3}\left(  1\right)  \beta\left(  1\right)  \otimes\varphi_{1}\left(
2\right)  \alpha\left(  2\right)  \right\} \nonumber\\
&  +\frac{1}{\sqrt{2}}\left\{  \varphi_{2}\left(  1\right)  \alpha\left(
1\right)  \otimes\varphi_{4}\left(  2\right)  \beta\left(  2\right)
-\varphi_{2}\left(  1\right)  \beta\left(  1\right)  \otimes\varphi_{4}\left(
2\right)  \alpha\left(  2\right)  \right\} \nonumber
\end{align}%
\begin{align}
&  =\frac{1}{\sqrt{2}}\left\{  \left\{  \varphi_{1}\left(  1\right)
\otimes\varphi_{3}\left(  2\right)  +\varphi_{2}\left(  1\right)
\otimes\varphi_{4}\left(  2\right)  \right\}  \alpha\left(  1\right)
\otimes\beta\left(  2\right)  \right. \nonumber\\
&  -\left\{  \varphi_{3}\left(  1\right)  \otimes\varphi_{1}\left(  2\right)
+\varphi_{4}\left(  1\right)  \otimes\varphi_{2}\left(  2\right)  \right\}
\beta\left(  1\right)  \otimes\alpha\left(  2\right)  .
\end{align}
One can observe that if we replace $\left\vert \psi_{2}\right\rangle $ with
$-\left\vert \psi_{2}\right\rangle $ $\equiv\left\vert \varphi_{2}%
\right\rangle \rightarrow-\left\vert \varphi_{2}\right\rangle $ we obtain
\[
\left\vert g_{-}\right\rangle =\left\vert \psi_{1}\psi_{3}\right\rangle
-\left\vert \psi_{2}\psi_{4}\right\rangle =\left\vert \varphi_{1}\alpha
\varphi_{3}\beta\right\rangle -\left\vert \varphi_{2}\alpha\varphi_{4}%
\beta\right\rangle
\]%
\begin{align}
&  =\frac{1}{\sqrt{2}}\left\{  \varphi_{1}\left(  1\right)  \alpha\left(
1\right)  \otimes\varphi_{3}\left(  2\right)  \beta\left(  2\right)
-\varphi_{3}\left(  1\right)  \beta\left(  1\right)  \otimes\varphi_{1}\left(
2\right)  \alpha\left(  2\right)  \right\} \nonumber\\
&  -\frac{1}{\sqrt{2}}\left\{  \varphi_{2}\left(  1\right)  \alpha\left(
1\right)  \otimes\varphi_{4}\left(  2\right)  \beta\left(  2\right)
-\varphi_{2}\left(  1\right)  \beta\left(  1\right)  \otimes\varphi_{4}\left(
2\right)  \alpha\left(  2\right)  \right\} \nonumber
\end{align}%
\begin{align}
&  =\frac{1}{\sqrt{2}}\left\{  \left\{  \varphi_{1}\left(  1\right)
\otimes\varphi_{3}\left(  2\right)  -\varphi_{2}\left(  1\right)
\otimes\varphi_{4}\left(  2\right)  \right\}  \alpha\left(  1\right)
\otimes\beta\left(  2\right)  \right. \nonumber\\
&  -\left\{  \varphi_{3}\left(  1\right)  \otimes\varphi_{1}\left(  2\right)
-\varphi_{2}\left(  1\right)  \otimes\varphi_{4}\left(  2\right)  \right\}
\beta\left(  1\right)  \otimes\alpha\left(  2\right)
\end{align}
and do not alter the FORDO, but do change the SORDO. The common FORDO\ is
\begin{align}
D^{1}  &  =\frac{1}{4}\left\{  \left\vert \psi_{1}\right\rangle \left\langle
\psi_{1}\right\vert +\left\vert \psi_{2}\right\rangle \left\langle \psi
_{2}\right\vert +\left\vert \psi_{3}\right\rangle \left\langle \psi
_{3}\right\vert +\left\vert \psi_{4}\right\rangle \left\langle \psi
_{4}\right\vert \right\} \nonumber\\
&  \equiv\frac{1}{4}\left\{  \left\vert 1\alpha\right\rangle \left\langle
1\alpha\right\vert +\left\vert 2\alpha\right\rangle \left\langle
2\alpha\right\vert +\left\vert 3\beta\right\rangle \left\langle 3\beta
\right\vert +\left\vert 4\beta\right\rangle \left\langle 4\beta\right\vert
\right\}
\end{align}
while the SORDO's are
\begin{align}
D^{2}\left(  g_{+}\right)   &  =\frac{1}{2}\left\{  \left\vert \psi_{1}%
\psi_{3}\right\rangle \left\langle \psi_{1}\psi_{3}\right\vert +\left\vert
\psi_{1}\psi_{3}\right\rangle \left\langle \psi_{2}\psi_{4}\right\vert
+\left\vert \psi_{2}\psi_{4}\right\rangle \left\langle \psi_{1}\psi
_{3}\right\vert +\left\vert \psi_{2}\psi_{4}\right\rangle \left\langle
\psi_{2}\psi_{4}\right\vert \right\} \nonumber\\
&  \equiv\frac{1}{2}\left\{  \left\vert 1\alpha3\beta\right\rangle
\left\langle 1\alpha3\beta\right\vert +\left\vert 1\alpha3\beta\right\rangle
\left\langle 2\alpha4\beta\right\vert +\left\vert 2\alpha4\beta\right\rangle
\left\langle 1\alpha3\beta\right\vert +\left\vert 2\alpha4\beta\right\rangle
\left\langle 2\alpha4\beta\right\vert \right\}  .
\end{align}
and
\begin{align}
D^{2}\left(  g_{-}\right)   &  =\frac{1}{2}\left\{  \left\vert \psi_{1}%
\psi_{3}\right\rangle \left\langle \psi_{1}\psi_{3}\right\vert -\left\vert
\psi_{1}\psi_{3}\right\rangle \left\langle \psi_{2}\psi_{4}\right\vert
-\left\vert \psi_{2}\psi_{4}\right\rangle \left\langle \psi_{1}\psi
_{3}\right\vert +\left\vert \psi_{2}\psi_{4}\right\rangle \left\langle
\psi_{2}\psi_{4}\right\vert \right\} \nonumber\\
&  \equiv\frac{1}{2}\left\{  \left\vert 1\alpha3\beta\right\rangle
\left\langle 1\alpha3\beta\right\vert -\left\vert 1\alpha3\beta\right\rangle
\left\langle 2\alpha4\beta\right\vert -\left\vert 2\alpha4\beta\right\rangle
\left\langle 1\alpha3\beta\right\vert +\left\vert 2\alpha4\beta\right\rangle
\left\langle 2\alpha4\beta\right\vert \right\}  .
\end{align}
If we further set
\begin{align}
\left\vert \varphi_{1}\right\rangle  &  =\left\vert \varphi_{3}\right\rangle
\nonumber\\
\left\vert \varphi_{2}\right\rangle  &  =\left\vert \varphi_{4}\right\rangle
\end{align}
to form restricted singlet spin geminals that factor into a product of space
and spin parts then
\begin{align}
\left\vert g_{+}\right\rangle  &  =\frac{1}{\sqrt{2}}\left\{  \overset{}%
{}\left\{  \varphi_{1}\left(  1\right)  \otimes\varphi_{1}\left(  2\right)
+\varphi_{2}\left(  1\right)  \otimes\varphi_{2}\left(  2\right)  \right\}
\right. \nonumber\\
&  \left.  \left\{  \alpha\left(  1\right)  \otimes\beta\left(  2\right)
-\beta\left(  1\right)  \otimes\alpha\left(  2\right)  \right\}  \overset{}%
{}\right\}
\end{align}
and
\begin{align}
\left\vert g_{-}\right\rangle  &  =\frac{1}{\sqrt{2}}\left\{  \overset{}%
{}\left\{  \varphi_{1}\left(  1\right)  \otimes\varphi_{1}\left(  2\right)
-\varphi_{2}\left(  1\right)  \otimes\varphi_{2}\left(  2\right)  \right\}
\right. \nonumber\\
&  \left.  \left\{  \alpha\left(  1\right)  \otimes\beta\left(  2\right)
-\beta\left(  1\right)  \otimes\alpha\left(  2\right)  \right\}  \overset{}%
{}\right\}
\end{align}
i.e.
\[
\left\vert g_{+}\right\rangle =\left(  \sqrt{2}\right)  ^{-1}\left\{
\left\vert \psi_{1}\psi_{3}\right\rangle +\left\vert \psi_{2}\psi
_{4}\right\rangle \right\}  =\left(  \sqrt{2}\right)  ^{-1}\left\{  \left\vert
\varphi_{1}\alpha\varphi_{1}\beta\right\rangle +\left\vert \varphi_{2}%
\alpha\varphi_{2}\beta\right\rangle \right\}
\]%
\[
\left\vert g_{-}\right\rangle =\left(  \sqrt{2}\right)  ^{-1}\left\{
\left\vert \psi_{1}\psi_{3}\right\rangle -\left\vert \psi_{2}\psi
_{4}\right\rangle \right\}  =\left(  \sqrt{2}\right)  ^{-1}\left\{  \left\vert
\varphi_{1}\alpha\varphi_{1}\beta\right\rangle -\left\vert \varphi_{2}%
\alpha\varphi_{2}\beta\right\rangle \right\}
\]
and $\left\vert g_{-}\right\rangle $ is obtained from $\left\vert
g_{+}\right\rangle $ by replacing $\left\vert \varphi_{2}\right\rangle $ by
$\left\vert i\varphi_{2}\right\rangle $. Thus
\[
D^{1}\left(  g_{+}\right)  =D^{1}\left(  g_{-}\right)
\]
while again
\[
D^{2}\left(  g_{+}\right)  \neq D^{2}\left(  g_{-}\right)
\]
as
\begin{align}
D^{2}\left(  g_{+}\right)   &  =\frac{1}{2}\left\{  \left\vert \psi_{1}%
\psi_{3}\right\rangle \left\langle \psi_{1}\psi_{3}\right\vert +\left\vert
\psi_{1}\psi_{3}\right\rangle \left\langle \psi_{2}\psi_{4}\right\vert
+\left\vert \psi_{2}\psi_{4}\right\rangle \left\langle \psi_{1}\psi
_{3}\right\vert +\left\vert \psi_{2}\psi_{4}\right\rangle \left\langle
\psi_{2}\psi_{4}\right\vert \right\} \nonumber\\
&  \equiv\frac{1}{2}\left\{  \left\vert 1\alpha1\beta\right\rangle
\left\langle 1\alpha1\beta\right\vert +\left\vert 1\alpha1\beta\right\rangle
\left\langle 2\alpha2\beta\right\vert +\left\vert 2\alpha2\beta\right\rangle
\left\langle 1\alpha1\beta\right\vert +\left\vert 2\alpha2\beta\right\rangle
\left\langle 2\alpha2\beta\right\vert \right\}  .
\end{align}
and
\begin{align}
D^{2}\left(  g_{-}\right)   &  =\frac{1}{2}\left\{  \left\vert \psi_{1}%
\psi_{3}\right\rangle \left\langle \psi_{1}\psi_{3}\right\vert -\left\vert
\psi_{1}\psi_{3}\right\rangle \left\langle \psi_{2}\psi_{4}\right\vert
-\left\vert \psi_{2}\psi_{4}\right\rangle \left\langle \psi_{1}\psi
_{3}\right\vert +\left\vert \psi_{2}\psi_{4}\right\rangle \left\langle
\psi_{2}\psi_{4}\right\vert \right\} \nonumber\\
&  \equiv\frac{1}{2}\left\{  \left\vert 1\alpha1\beta\right\rangle
\left\langle 1\alpha1\beta\right\vert -\left\vert 1\alpha1\beta\right\rangle
\left\langle 2\alpha2\beta\right\vert -\left\vert 2\alpha2\beta\right\rangle
\left\langle 1\alpha1\beta\right\vert +\left\vert 2\alpha2\beta\right\rangle
\left\langle 2\alpha2\beta\right\vert \right\}  .
\end{align}


\subsection{SORDO}

We will use the general form of the SORDO for the case
\begin{align*}
m  &  =\frac{N}{2}=\frac{2}{2}=1\\
n_{1}  &  =n_{2}=\frac{1}{2}\\
c_{1}  &  =c_{2}=\frac{1}{\sqrt{2}}.
\end{align*}
As $m=1$ the symmetric functions, which refer to the number of pairs are
\begin{align*}
S_{1-1}\left(  \hat{\jmath}\hat{k}\right)   &  =S_{0}\left(  \hat{\jmath}%
\hat{k}\right)  =1\\
S_{1-1}\left(  \hat{\jmath}\right)   &  =S_{0}\left(  \hat{\jmath}\right)
=1\\
S_{1-2}\left(  \hat{\jmath}\hat{k}\right)   &  =S_{-1}\left(  \hat{\jmath}%
\hat{k}\right)  =0.
\end{align*}
Using the general form we obtain%
\begin{align*}
&  D^{2}\left(  \left\vert g^{1}\right\rangle \left\langle g^{1}\right\vert
\right)  =\left(  S_{1}\right)  ^{-1}\left\{  \sum_{1\leq j\leq s}n_{j}%
S_{1-1}\left(  \hat{\jmath}\right)  \left\vert \psi_{j}\psi_{j+s}\right\rangle
\left\langle \psi_{j}\psi_{j+s}\right\vert \right. \\
&  +\sum_{1\leq j\neq k\leq s}c_{j}c_{k}e^{i\left(  \theta_{j}-\theta
_{k}\right)  }S_{1-1}\left(  \hat{\jmath}\hat{k}\right)  \left\vert \psi
_{j}\psi_{j+s}\right\rangle \left\langle \psi_{k}\psi_{k+s}\right\vert \left.
+\sum_{\substack{1\leq j<k\leq r\\k\neq j+s}}n_{j}n_{k}S_{1-2}\left(
\hat{\jmath}\hat{k}\right)  \left\vert \psi_{j}\psi_{k}\right\rangle
\left\langle \psi_{j}\psi_{k}\right\vert \right\}  ,
\end{align*}
which gives, after noting
\begin{align*}
S_{1}  &  =n_{1}+n_{2}=1\\
S_{\frac{2}{2}-1}\left(  \hat{1}\hat{2}\right)   &  =S_{0}\left(  \hat{1}%
\hat{2}\right)  =1\\
S_{\frac{2}{2}-1}\left(  \hat{1}\right)   &  =S_{0}\left(  \hat{1}\right)
=1;S_{\frac{2}{2}-1}\left(  \hat{2}\right)  =S_{0}\left(  \hat{2}\right)  =1\\
S_{\frac{2}{2}-2}\left(  \hat{1}\hat{2}\right)   &  =S_{-1}\left(  \hat{1}%
\hat{2}\right)  =0,
\end{align*}
the SORDO%
\begin{align*}
&  D^{2}\left(  \left\vert g_{+}^{1}\right\rangle \left\langle g_{+}%
^{1}\right\vert \right)  =\frac{1}{2}\left\{  \left\vert \psi_{1}\psi
_{3}\right\rangle \left\langle \psi_{1}\psi_{3}\right\vert +\left\vert
\psi_{2}\psi_{4}\right\rangle \left\langle \psi_{2}\psi_{4}\right\vert
\right\} \\
&  +\frac{1}{2}\left\{  \left\vert \psi_{1}\psi_{3}\right\rangle \left\langle
\psi_{2}\psi_{4}\right\vert e^{i\left(  \theta_{1}-\theta_{2}\right)
}+\left\vert \psi_{2}\psi_{4}\right\rangle \left\langle \psi_{1}\psi
_{3}\right\vert e^{-i\left(  \theta_{1}-\theta_{2}\right)  }\right\} \\
&  =\frac{1}{2}\left\{  \left\vert \psi_{1}\psi_{3}\right\rangle \left\langle
\psi_{1}\psi_{3}\right\vert +\left\vert \psi_{1}\psi_{3}\right\rangle
\left\langle \psi_{2}\psi_{4}\right\vert +\left\vert \psi_{2}\psi
_{4}\right\rangle \left\langle \psi_{1}\psi_{3}\right\vert +\left\vert
\psi_{2}\psi_{4}\right\rangle \left\langle \psi_{2}\psi_{4}\right\vert
\right\}
\end{align*}
as
\[
\theta_{1}=\theta_{2}=0.
\]
This should be compared to the previous direct derivation%
\begin{align}
D^{2}\left(  g_{+}\right)   &  =\frac{1}{2}\left\{  \left\vert \psi_{1}%
\psi_{3}\right\rangle \left\langle \psi_{1}\psi_{3}\right\vert +\left\vert
\psi_{1}\psi_{3}\right\rangle \left\langle \psi_{2}\psi_{4}\right\vert
+\left\vert \psi_{2}\psi_{4}\right\rangle \left\langle \psi_{1}\psi
_{3}\right\vert +\left\vert \psi_{2}\psi_{4}\right\rangle \left\langle
\psi_{2}\psi_{4}\right\vert \right\} \nonumber\\
&  \equiv\frac{1}{2}\left\{  \left\vert 1\alpha3\beta\right\rangle
\left\langle 1\alpha3\beta\right\vert +\left\vert 1\alpha3\beta\right\rangle
\left\langle 2\alpha4\beta\right\vert +\left\vert 2\alpha4\beta\right\rangle
\left\langle 1\alpha3\beta\right\vert +\left\vert 2\alpha4\beta\right\rangle
\left\langle 2\alpha4\beta\right\vert \right\}  .
\end{align}
In the $\left\vert g_{-}\right\rangle $ case
\begin{align*}
\theta_{1}  &  =0\\
\theta_{2}  &  =\pi
\end{align*}
and we obtain
\begin{align*}
&  D^{2}\left(  \left\vert g_{-}^{1}\right\rangle \left\langle g_{-}%
^{1}\right\vert \right)  =\frac{1}{2}\left\{  \left\vert \psi_{1}\psi
_{3}\right\rangle \left\langle \psi_{1}\psi_{3}\right\vert +\left\vert
\psi_{2}\psi_{4}\right\rangle \left\langle \psi_{2}\psi_{4}\right\vert
\right\} \\
&  +\frac{1}{2}\left\{  \left\vert \psi_{1}\psi_{3}\right\rangle \left\langle
\psi_{2}\psi_{4}\right\vert e^{i\left(  \theta_{1}-\theta_{2}\right)
}+\left\vert \psi_{2}\psi_{4}\right\rangle \left\langle \psi_{1}\psi
_{3}\right\vert e^{-i\left(  \theta_{1}-\theta_{2}\right)  }\right\} \\
&  =\frac{1}{2}\left\{  \left\vert \psi_{1}\psi_{3}\right\rangle \left\langle
\psi_{1}\psi_{3}\right\vert -\left\vert \psi_{1}\psi_{3}\right\rangle
\left\langle \psi_{2}\psi_{4}\right\vert -\left\vert \psi_{2}\psi
_{4}\right\rangle \left\langle \psi_{1}\psi_{3}\right\vert +\left\vert
\psi_{2}\psi_{4}\right\rangle \left\langle \psi_{2}\psi_{4}\right\vert
\right\}
\end{align*}
as
\begin{align*}
e^{i\theta}  &  =\cos\theta-i\sin\theta\\
e^{i\pi}  &  =\cos\pi-i\sin\pi=-1\\
e^{-i\pi}  &  =\cos\left(  -\pi\right)  -i\sin\left(  -\pi\right)  =-1.
\end{align*}


\subsection{Energy}

In these cases
\[
\left\vert g_{\theta}\right\rangle =c_{1}e^{i\theta_{1}}\left\vert \psi
_{1}\psi_{3}\right\rangle +c_{2}e^{i\theta_{2}}\left\vert \psi_{2}\psi
_{4}\right\rangle
\]


and%
\begin{align*}
n_{1}  &  =n_{2}=\frac{1}{2}\\
c_{1}  &  =c_{2}=\frac{1}{\sqrt{2}}%
\end{align*}
giving%
\begin{align*}
\left\vert g_{\theta}\right\rangle  &  =\frac{1}{\sqrt{2}}\left\{  \left\vert
\psi_{1}\psi_{3}\right\rangle +e^{i\theta}\left\vert \psi_{2}\psi
_{4}\right\rangle \right\} \\
&  =\frac{1}{\sqrt{2}}\left\{  \left\vert 1\alpha1\beta\right\rangle
+e^{i\theta}\left\vert 2\alpha2\beta\right\rangle \right\}  .
\end{align*}
The symmetric functions are
\begin{align*}
S_{1}\left(  \mathbf{n}\right)   &  =n_{1}+n_{2}=1\\
S_{0}\left(  \hat{1}\right)   &  =S_{0}\left(  \hat{2}\right)  =1\\
S_{-1}\left(  \hat{1}\hat{2}\right)   &  =0
\end{align*}
and the energy is given by%
\begin{align*}
&  \mathcal{E}\left(  \theta\right)  =\frac{\mathcal{H}\left(  \theta\right)
}{S_{\frac{N}{2}}}=S_{1}\left(  \mathbf{n}\right)  ^{-1}\left\{  \sum_{1\leq
j\leq2}n_{j}S_{0}\left(  \hat{\jmath}\right)  \left\{  \overset{}{\left\langle
\psi_{j}\left\vert h_{1}\psi_{j}\right.  \right\rangle +\left\langle
\psi_{j+s}\left\vert h_{1}\psi_{j+s}\right.  \right\rangle }+\,\left\langle
\psi_{j}\psi_{j+s}\right\vert \left\vert \psi_{j}\psi_{j+s}\right\rangle
\right\}  \right. \\
&  +\sum_{1\leq j\neq k\leq2}\left\{  c_{j}c_{k}e^{i\left(  \theta_{j}%
-\theta_{k}\right)  }S_{0}\left(  \hat{\jmath}\hat{k}\right)  \left\langle
\psi_{j}\psi_{j+s}\right\vert \left\vert \psi_{k}\psi_{k+s}\right\rangle
+\frac{1}{2}\overset{}{n_{i}n_{j}S_{-1}\left(  \hat{\jmath}\hat{k}\right)
\left\langle \psi_{j}\psi_{k}|||\psi_{j}\psi_{k}\right\rangle }\right\}
\end{align*}%
\begin{align*}
&  =\frac{1}{2}\sum_{1\leq j\leq2}\left\{  \overset{}{\left\langle \psi
_{j}\left\vert h_{1}\psi_{j}\right.  \right\rangle +\left\langle \psi
_{j+s}\left\vert h_{1}\psi_{j+s}\right.  \right\rangle }+\,\left\langle
\psi_{j}\psi_{j+s}\right\vert \left\vert \psi_{j}\psi_{j+s}\right\rangle
\right\} \\
&  +\frac{1}{2}\sum_{1\leq j\neq k\leq2}e^{i\left(  \theta_{j}-\theta
_{k}\right)  }\left\langle \psi_{j}\psi_{j+s}\right\vert \left\vert \psi
_{k}\psi_{k+s}\right\rangle \\
&  =\frac{1}{2}\left\{  \overset{}{\left\langle \psi_{1}\left\vert h_{1}%
\psi_{1}\right.  \right\rangle +\left\langle \psi_{2}\left\vert h_{1}\psi
_{2}\right.  \right\rangle +\left\langle \psi_{3}\left\vert h_{1}\psi
_{3}\right.  \right\rangle +\left\langle \psi_{4}\left\vert h_{1}\psi
_{4}\right.  \right\rangle }\right. \\
&  \left.  \,+\left\langle \psi_{1}\psi_{3}\right\vert \left\vert \psi_{1}%
\psi_{3}\right\rangle +\left\langle \psi_{2}\psi_{4}\right\vert \left\vert
\psi_{2}\psi_{4}\right\rangle +e^{i\theta}\left\langle \psi_{1}\psi
_{3}\right\vert \left\vert \psi_{2}\psi_{4}\right\rangle +e^{-i\theta
}\left\langle \psi_{2}\psi_{4}\right\vert \left\vert \psi_{1}\psi
_{3}\right\rangle \overset{}{}\right\}  .
\end{align*}
Thus giving in the +/- cases%
\begin{align*}
&  \mathcal{E}\left(  g_{+}\right)  =\frac{1}{2}\left\{  \overset
{}{\left\langle 1\alpha\left\vert h_{1}1\alpha\right.  \right\rangle
+\left\langle 2\alpha\left\vert h_{1}2\alpha\right.  \right\rangle
+\left\langle 1\beta\left\vert h_{1}1\beta\right.  \right\rangle +\left\langle
2\beta\left\vert h_{1}2\beta\right.  \right\rangle }\right. \\
&  \left.  \,+\left\langle 1\alpha1\beta\right\vert \left\vert 1\alpha
1\beta\right\rangle +\left\langle 2\alpha2\beta\right\vert \left\vert
2\alpha2\beta\right\rangle +\left\langle 1\alpha1\beta\right\vert \left\vert
2\alpha2\beta\right\rangle +\left\langle 1\alpha1\beta\right\vert \left\vert
2\alpha2\beta\right\rangle \overset{}{}\right\}
\end{align*}
and
\begin{align*}
&  \mathcal{E}\left(  g_{-}\right)  =\frac{1}{2}\left\{  \overset
{}{\left\langle 1\alpha\left\vert h_{1}1\alpha\right.  \right\rangle
+\left\langle 2\alpha\left\vert h_{1}2\alpha\right.  \right\rangle
+\left\langle 1\beta\left\vert h_{1}1\beta\right.  \right\rangle +\left\langle
2\beta\left\vert h_{1}2\beta\right.  \right\rangle }\right. \\
&  \left.  \,+\left\langle 1\alpha1\beta\right\vert \left\vert 1\alpha
1\beta\right\rangle +\left\langle 2\alpha2\beta\right\vert \left\vert
2\alpha2\beta\right\rangle -\left\langle 1\alpha1\beta\right\vert \left\vert
2\alpha2\beta\right\rangle -\left\langle 1\alpha1\beta\right\vert \left\vert
2\alpha2\beta\right\rangle \overset{}{}\right\}  ,
\end{align*}
where%
\begin{align*}
&  \left\langle \psi_{j}\psi_{k}|||\psi_{j}\psi_{k}\right\rangle =\left\langle
jk||jk\right\rangle +\left\langle j\,k+s||j\,k+s\right\rangle +\\
&  \left\langle j+s\,k||j+s\,k\right\rangle +\left\langle
j+s\,k+s||j+s\,k+s\right\rangle ;1\leq j,k\leq s.
\end{align*}
and%

\begin{align*}
\left\langle jk||lm\right\rangle  &  =\int\left\{  \psi_{j}^{\ast}\left(
\mathbf{r}_{1}\mathbf{,\xi}_{1}\right)  \psi_{k}^{\ast}\left(  \mathbf{r}%
_{2}\mathbf{,\xi}_{2}\right)  \frac{1}{\left\Vert \mathbf{r}_{1}%
-\mathbf{r}_{2}\right\Vert }\psi_{l}\left(  \mathbf{r}_{1}\mathbf{,\xi}%
_{1}\right)  \psi_{m}\left(  \mathbf{r}_{2}\mathbf{,\xi}_{2}\right)  \right.
\\
&  -\left.  \psi_{j}^{\ast}\left(  \mathbf{r}_{1}\mathbf{,\xi}_{1}\right)
\psi_{k}^{\ast}\left(  \mathbf{r}_{2}\mathbf{,\xi}_{2}\right)  \frac
{1}{\left\Vert \mathbf{r}_{1}-\mathbf{r}_{2}\right\Vert }\psi_{l}\left(
\mathbf{r}_{2}\mathbf{,\xi}_{2}\right)  \psi_{m}\left(  \mathbf{r}%
_{1}\mathbf{,\xi}_{1}\right)  \right\}  d\mathbf{r}_{1}d\mathbf{r}_{2}d\xi
_{1}d\xi_{2}\\
&  \left.  {}\right.  \qquad\qquad\qquad\qquad\qquad\qquad\qquad\qquad
\qquad1\leq j,k\leq2s;1\leq l,m\leq2s.
\end{align*}


\section{Pure Spin Geminals}

We can compare the preceding $\left\vert g_{+}\right\rangle $ and $\left\vert
g_{-}\right\rangle $ geminals with pure spin ones that do not have the same
FORDO. The low spin triplet is
\[
\left\vert g_{10}\right\rangle =\left(  \sqrt{2}\right)  ^{-1}\left\{
\left\vert \psi_{1}\psi_{3}\right\rangle +\left\vert \psi_{2}\psi
_{4}\right\rangle \right\}  =\left(  \sqrt{2}\right)  ^{-1}\left\{  \left\vert
\varphi_{1}\alpha\varphi_{2}\beta\right\rangle +\left\vert \varphi_{1}%
\beta\varphi_{2}\alpha\right\rangle \right\}  ,
\]
which can be seperated into a product of spacr and spin factors as%
\begin{align}
\left\vert g_{10}\right\rangle  &  =\left(  \sqrt{2}\right)  ^{-1}\left\{
\left\vert \varphi_{1}\alpha\varphi_{2}\beta\right\rangle +\left\vert
\varphi_{1}\beta\varphi_{2}\alpha\right\rangle \right\} \nonumber\\
&  =\frac{1}{\sqrt{2}}\left\{  \varphi_{1}\left(  1\right)  \alpha\left(
1\right)  \otimes\varphi_{2}\left(  2\right)  \beta\left(  2\right)
-\varphi_{2}\left(  1\right)  \beta\left(  1\right)  \otimes\varphi_{1}\left(
2\right)  \alpha\left(  2\right)  \right\} \nonumber\\
&  +\frac{1}{\sqrt{2}}\left\{  \varphi_{1}\left(  1\right)  \beta\left(
1\right)  \otimes\varphi_{2}\left(  2\right)  \alpha\left(  2\right)
-\varphi_{2}\left(  1\right)  \alpha\left(  1\right)  \otimes\varphi
_{1}\left(  2\right)  \beta\left(  2\right)  \right\} \nonumber\\
&  =\frac{1}{\sqrt{2}}\left\{  \left\{  \varphi_{1}\left(  1\right)
\otimes\varphi_{2}\left(  2\right)  -\varphi_{2}\left(  1\right)
\otimes\varphi_{1}\left(  2\right)  \right\}  \alpha\left(  1\right)
\otimes\beta\left(  2\right)  \right. \nonumber\\
&  +\left\{  \varphi_{2}\left(  1\right)  \otimes\varphi_{1}\left(  2\right)
-\varphi_{1}\left(  1\right)  \otimes\varphi_{2}\left(  2\right)  \right\}
\beta\left(  1\right)  \otimes\alpha\left(  2\right) \nonumber\\
&  =\frac{1}{\sqrt{2}}\left\{  \left\{  \varphi_{1}\left(  1\right)
\otimes\varphi_{2}\left(  2\right)  -\varphi_{2}\left(  1\right)
\otimes\varphi_{1}\left(  2\right)  \right\}  \left\{  \alpha\left(  1\right)
\otimes\beta\left(  2\right)  +\beta\left(  1\right)  \otimes\alpha\left(
2\right)  \right\}  \right\}  ,
\end{align}
and the low spin singlet%
\[
\left\vert g_{00}\right\rangle =\left(  \sqrt{2}\right)  ^{-1}\left\{
\left\vert \psi_{1}\psi_{3}\right\rangle -\left\vert \psi_{2}\psi
_{4}\right\rangle \right\}  =\left(  \sqrt{2}\right)  ^{-1}\left\{  \left\vert
\varphi_{1}\alpha\varphi_{2}\beta\right\rangle -\left\vert \varphi_{1}%
\beta\varphi_{2}\alpha\right\rangle \right\}  ,
\]
which can be seperated as%
\begin{align}
\left\vert g_{00}\right\rangle  &  =\left(  \sqrt{2}\right)  ^{-1}\left\{
\left\vert \varphi_{1}\alpha\varphi_{2}\beta\right\rangle -\left\vert
\varphi_{1}\beta\varphi_{2}\alpha\right\rangle \right\} \nonumber\\
&  =\frac{1}{\sqrt{2}}\left\{  \varphi_{1}\left(  1\right)  \alpha\left(
1\right)  \otimes\varphi_{2}\left(  2\right)  \beta\left(  2\right)
-\varphi_{2}\left(  1\right)  \beta\left(  1\right)  \otimes\varphi_{1}\left(
2\right)  \alpha\left(  2\right)  \right\} \nonumber\\
&  -\frac{1}{\sqrt{2}}\left\{  \varphi_{1}\left(  1\right)  \beta\left(
1\right)  \otimes\varphi_{2}\left(  2\right)  \alpha\left(  2\right)
-\varphi_{2}\left(  1\right)  \alpha\left(  1\right)  \otimes\varphi
_{1}\left(  2\right)  \beta\left(  2\right)  \right\} \nonumber\\
&  =\frac{1}{\sqrt{2}}\left\{  \left\{  \varphi_{1}\left(  1\right)
\otimes\varphi_{2}\left(  2\right)  +\varphi_{2}\left(  1\right)
\otimes\varphi_{1}\left(  2\right)  \right\}  \alpha\left(  1\right)
\otimes\beta\left(  2\right)  \right. \nonumber\\
&  -\left\{  \varphi_{2}\left(  1\right)  \otimes\varphi_{1}\left(  2\right)
+\varphi_{1}\left(  1\right)  \otimes\varphi_{2}\left(  2\right)  \right\}
\beta\left(  1\right)  \otimes\alpha\left(  2\right) \nonumber\\
&  =\frac{1}{\sqrt{2}}\left\{  \left\{  \varphi_{1}\left(  1\right)
\otimes\varphi_{2}\left(  2\right)  +\varphi_{2}\left(  1\right)
\otimes\varphi_{1}\left(  2\right)  \right\}  \left\{  \alpha\left(  1\right)
\otimes\beta\left(  2\right)  -\beta\left(  1\right)  \otimes\alpha\left(
2\right)  \right\}  \right\}  .
\end{align}


The spin squared operator is given by
\[
S^{2}=S_{-}S_{+}+S_{0}^{2}+S_{0}%
\]
where
\begin{align*}
S_{+}  &  =\sum_{1\leq j\leq s}a_{j\alpha}^{\dagger}a_{j\beta}\\
S_{-}  &  =\sum_{1\leq j\leq s}a_{j\beta}^{\dagger}a_{j\alpha}=S_{+}^{\dagger
}\\
S_{0}  &  =\sum_{1\leq j\leq s}a_{j\alpha}^{\dagger}a_{j\alpha}-\sum_{1\leq
j\leq s}a_{j\beta}^{\dagger}a_{j\beta}.
\end{align*}
Thus
\begin{align*}
&  S^{2}\left\vert g_{10}\right\rangle =S^{2}\left(  \sqrt{2}\right)
^{-1}\left\{  \left\vert \varphi_{1}\alpha\varphi_{2}\beta\right\rangle
+\left\vert \varphi_{1}\beta\varphi_{2}\alpha\right\rangle \right\}
=S_{-}S_{+}\left(  \sqrt{2}\right)  ^{-1}\left\{  \left\vert \varphi_{1}%
\alpha\varphi_{2}\beta\right\rangle +\left\vert \varphi_{1}\beta\varphi
_{2}\alpha\right\rangle \right\} \\
&  =S_{-}\left(  \sqrt{2}\right)  ^{-1}\left\{  \left\vert \varphi_{1}%
\alpha\varphi_{2}\alpha\right\rangle +\left\vert \varphi_{1}\alpha\varphi
_{2}\alpha\right\rangle \right\}  =2\left(  \sqrt{2}\right)  ^{-1}\left\{
\left\vert \varphi_{1}\beta\varphi_{2}\alpha\right\rangle +\left\vert
\varphi_{1}\alpha\varphi_{2}\beta\right\rangle \right\}  =2\left\vert
g_{10}\right\rangle
\end{align*}
and
\begin{align*}
&  S^{2}\left\vert g_{00}\right\rangle =S^{2}\left(  \sqrt{2}\right)
^{-1}\left\{  \left\vert \varphi_{1}\alpha\varphi_{2}\beta\right\rangle
-\left\vert \varphi_{1}\beta\varphi_{2}\alpha\right\rangle \right\}
=S_{-}S_{+}\left(  \sqrt{2}\right)  ^{-1}\left\{  \left\vert \varphi_{1}%
\alpha\varphi_{2}\beta\right\rangle -\left\vert \varphi_{1}\beta\varphi
_{2}\alpha\right\rangle \right\} \\
&  =S_{-}\left(  \sqrt{2}\right)  ^{-1}\left\{  \left\vert \varphi_{1}%
\alpha\varphi_{2}\alpha\right\rangle -\left\vert \varphi_{1}\alpha\varphi
_{2}\alpha\right\rangle \right\} \\
&  =\left(  \sqrt{2}\right)  ^{-1}\left\{  \left\vert \varphi_{1}\beta
\varphi_{2}\alpha\right\rangle +\left\vert \varphi_{1}\alpha\varphi_{2}%
\beta\right\rangle -\left\vert \varphi_{1}\beta\varphi_{2}\alpha\right\rangle
-\left\vert \varphi_{1}\beta\varphi_{2}\alpha\right\rangle \right\}  =0,
\end{align*}
showing that $\left\vert g_{10}\right\rangle $ and $\left\vert g_{00}%
\right\rangle $ are pure spin states.

This should be compared to%
\begin{align*}
S^{2}\left\vert g_{+}\right\rangle =  &  S^{2}\left(  \sqrt{2}\right)
^{-1}\left\{  \left\vert \varphi_{1}\alpha\varphi_{3}\beta\right\rangle
+\left\vert \varphi_{2}\alpha\varphi_{4}\beta\right\rangle \right\}
=S_{-}S_{+}\left(  \sqrt{2}\right)  ^{-1}\left\{  \left\vert \varphi_{1}%
\alpha\varphi_{3}\beta\right\rangle +\left\vert \varphi_{2}\alpha\varphi
_{4}\beta\right\rangle \right\} \\
&  =S_{-}\left(  \sqrt{2}\right)  ^{-1}\left\{  \left\vert \varphi_{1}%
\alpha\varphi_{3}\alpha\right\rangle +\left\vert \varphi_{2}\alpha\varphi
_{4}\alpha\right\rangle \right\} \\
&  =\left(  \sqrt{2}\right)  ^{-1}\left\{  \left\vert \varphi_{1}\beta
\varphi_{3}\alpha\right\rangle +\left\vert \varphi_{1}\alpha\varphi_{3}%
\beta\right\rangle +\left\vert \varphi_{2}\beta\varphi_{4}\alpha\right\rangle
+\left\vert \varphi_{2}\alpha\varphi_{4}\beta\right\rangle \right\}
\end{align*}
which shows that $\left\vert g_{+}\right\rangle $ is not in general an
eigenstate of $S^{2}$. However in the special case restricted case when
$\varphi_{1}=\varphi_{3}$ we obtain
\[
S^{2}\left\vert g_{+}\right\rangle =\left\vert 1\beta1\alpha\right\rangle
+\left\vert 1\alpha1\beta\right\rangle +\left\vert 2\beta2\alpha\right\rangle
+\left\vert 2\alpha2\beta\right\rangle =0,
\]
showing that it is a pure singlet state. The negative geminal gives
\begin{align*}
S^{2}\left\vert g_{-}\right\rangle =  &  S^{2}\left(  \sqrt{2}\right)
^{-1}\left\{  \left\vert \varphi_{1}\alpha\varphi_{3}\beta\right\rangle
-\left\vert \varphi_{2}\alpha\varphi_{4}\beta\right\rangle \right\}
=S_{-}S_{+}\left(  \sqrt{2}\right)  ^{-1}\left\{  \left\vert \varphi_{1}%
\alpha\varphi_{3}\beta\right\rangle -\left\vert \varphi_{2}\alpha\varphi
_{4}\beta\right\rangle \right\} \\
&  =S_{-}\left(  \sqrt{2}\right)  ^{-1}\left\{  \left\vert \varphi_{1}%
\alpha\varphi_{3}\alpha\right\rangle -\left\vert \varphi_{2}\alpha\varphi
_{4}\alpha\right\rangle \right\} \\
&  =\left(  \sqrt{2}\right)  ^{-1}\left\{  \left\vert \varphi_{1}\beta
\varphi_{3}\alpha\right\rangle +\left\vert \varphi_{1}\alpha\varphi_{3}%
\beta\right\rangle -\left\vert \varphi_{2}\beta\varphi_{4}\alpha\right\rangle
-\left\vert \varphi_{2}\alpha\varphi_{4}\beta\right\rangle \right\}  ,
\end{align*}
which again is not in general an eigenstate of spin unless $\varphi
_{1}=\varphi_{3}$ where again it becames a restricted singlet.

\section{Density Matrix Functional}

One can define an AGP DMF by
\[
F\left(  D^{1}\right)  =\min_{\mathbf{\theta}}\mathcal{E}\left(
\mathbf{\theta}\right)
\]
where for fixed $\mathbf{c}$ and $\mathbf{\psi}$ the energy is just a function
of the angular differences $\left\{  \left.  \theta_{j}-\theta_{k}\right\vert
1\leq j<k\leq s\right\}  $ and is given by
\begin{align*}
&  \mathcal{E}\left(  \mathbf{\theta}\right)  =\frac{\mathcal{H}\left(
\theta\right)  }{S_{\frac{N}{2}}}=S_{\frac{N}{2}}{}^{-1}\left\{  \sum_{1\leq
j\leq s}n_{j}S_{\frac{N}{2}-1}\left(  \hat{\jmath}\right)  \left\{  \overset
{}{\left\langle \psi_{j}\left\vert h_{1}\psi_{j}\right.  \right\rangle
+\left\langle \psi_{j+s}\left\vert h_{1}\psi_{j+s}\right.  \right\rangle
}+\,\left\langle \psi_{j}\psi_{j+s}\right\vert \left\vert \psi_{j}\psi
_{j+s}\right\rangle \right\}  \right. \\
&  +\sum_{1\leq j\neq k\leq s}\left\{  c_{j}c_{k}e^{i\left(  \theta_{j}%
-\theta_{k}\right)  }S_{\frac{N}{2}-1}\left(  \hat{\jmath}\hat{k}\right)
\left\langle \psi_{j}\psi_{j+s}\right\vert \left\vert \psi_{k}\psi
_{k+s}\right\rangle +\frac{1}{2}n_{j}n_{k}S_{\frac{N}{2}-2}\left(  \hat
{\jmath}\hat{k}\right)  \left\langle \psi_{j}\psi_{k}|||\psi_{j}\psi
_{k}\right\rangle \,\overset{}{}\right\}  .
\end{align*}
When the FORDO\ is doubly degenerate it is uniquely parameterized by $\left\{
\mathbf{\psi,c}\right\}  ,$ where $\mathbf{\psi}$ are representatives of the
equivalence classes produced by the normalized GSO's $\left\{  \left\vert
\psi\right\rangle \right\}  $ defined by
\begin{align*}
\left\{  \left\vert \psi_{a}\right\rangle ,\left\vert \psi_{b}\right\rangle
\right\}   &  \sim\left\{  \left\vert \psi_{c}\right\rangle ,\left\vert
\psi_{d}\right\rangle \right\} \\
&  \Updownarrow\\
\left\vert \psi_{a}\psi_{b}\right\rangle \left\langle \psi_{a}\psi
_{b}\right\vert  &  =\left\vert \psi_{c}\psi_{d}\right\rangle \left\langle
\psi_{c}\psi_{d}\right\vert
\end{align*}
the canonical coeffcients $\mathbf{c}$ lie on the positive portion of the
hypersphere
\[
\sum\limits_{1\leq j\leq s}c_{j}^{2}=1;\,c_{j}\geq0;1\leq j\leq s
\]
and the angles $\mathbf{\theta}=\left(  \theta_{2},\cdots,\theta_{s}\right)  $ satisfy%

\[
0\leq\theta_{2},\cdots,\theta_{s}<2\pi;\theta_{1}=0.
\]
In the case of our two particle example
\begin{align*}
\mathcal{E}\left(  \theta\right)   &  =\frac{1}{2}\left\{  \overset
{}{\left\langle 1\alpha\left\vert h_{1}1\alpha\right.  \right\rangle
+\left\langle 1\beta\left\vert h_{1}1\beta\right.  \right\rangle +\left\langle
2\alpha\left\vert h_{1}2\alpha\right.  \right\rangle +\left\langle
2\beta\left\vert h_{1}2\beta\right.  \right\rangle }\right.  \\
&  \left.  \,+\left\langle 1\alpha1\beta\right\vert \left\vert 1\alpha
1\beta\right\rangle +\left\langle 2\alpha2\beta\right\vert \left\vert
2\alpha2\beta\right\rangle +e^{i\theta}\left\langle 1\alpha1\beta\right\vert
\left\vert 2\alpha2\beta\right\rangle +e^{-i\theta}\left\langle 2\alpha
2\beta\right\vert \left\vert 1\alpha1\beta\right\rangle \overset{}{}\right\}
\\
&  =\frac{1}{2}\left\{  \overset{}{\left\langle 1\alpha\left\vert h_{1}%
1\alpha\right.  \right\rangle +\left\langle 1\beta\left\vert h_{1}%
1\beta\right.  \right\rangle +\left\langle 2\alpha\left\vert h_{1}%
2\alpha\right.  \right\rangle +\left\langle 2\beta\left\vert h_{1}%
2\beta\right.  \right\rangle }\right.  \\
&  \left.  \,+\left\langle 1\alpha1\beta\right\vert \left\vert 1\alpha
1\beta\right\rangle +\left\langle 2\alpha2\beta\right\vert \left\vert
2\alpha2\beta\right\rangle +2\operatorname{Re}e^{i\theta}\left\langle
1\alpha1\beta\right\vert \left\vert 2\alpha2\beta\right\rangle \overset{}%
{}\right\}
\end{align*}
which becomes%
\[
\mathcal{E}\left(  \theta\right)  =\left\langle 1\left\vert h_{1}1\right.
\right\rangle +\left\langle 2\left\vert h_{1}2\right.  \right\rangle +\left(
\left.  11\right\vert 11\right)  +\left(  \left.  22\right\vert 22\right)
+2\operatorname{Re}e^{i\theta}\left(  \left.  11\right\vert 22\right)  .
\]
The constraint on the canonical coefficients gives
\[
c_{1}^{2}+c_{2}^{2}=1
\]
leading to
\[
c_{2}=+\sqrt{1-c_{1}^{2}}%
\]
or the parameterization
\begin{align*}
c_{1} &  =\cos\chi\\
c_{2} &  =\sin\chi\\
0 &  \leq\chi\leq\frac{\pi}{2}.
\end{align*}
The specific value of the DMF at the point
\[
\left(  \mathbf{\varphi,c}\right)  =\left(  \varphi_{1},\varphi_{2},\frac
{1}{\sqrt{2}},\frac{1}{\sqrt{2}}\right)
\]
as
\[
F\left(  \mathbf{\varphi,}\left(  \frac{1}{\sqrt{2}},\frac{1}{\sqrt{2}%
}\right)  \right)  =\min_{\mathbf{\theta}}\left\{  \left\langle 1\left\vert
h_{1}1\right.  \right\rangle +\left\langle 2\left\vert h_{1}2\right.
\right\rangle +\left(  \left.  11\right\vert 11\right)  +\left(  \left.
22\right\vert 22\right)  +2\operatorname{Re}\left\{  e^{i\theta}\left(
\left.  11\right\vert 22\right)  \right\}  \right\}
\]
Note that
\begin{align*}
z  & =re^{i\omega}=r\left(  \cos\omega+i\sin\omega\right)  \\
r  & =+\sqrt{\left(  \operatorname{Re}z\right)  ^{2}+\left(  \operatorname{Im}%
z\right)  ^{2}}\\
\operatorname{Re}z  & =r\cos\omega;\,\operatorname{Im}z=r\sin\omega
\end{align*}
hence
\[
\operatorname{Re}\left\{  e^{i\theta}\left(  \left.  11\right\vert 22\right)
\right\}  =\operatorname{Re}\left\{  e^{i\theta}ae^{i\alpha}\right\}
=\operatorname{Re}\left\{  e^{i\left(  \theta+\alpha\right)  }a\right\}
=a\cos\left(  \theta+\alpha\right)
\]
where
\[
a=\left\vert \left(  \left.  11\right\vert 22\right)  \right\vert
\]
and thus
\begin{align*}
\min_{\theta}\left\{  a\cos\left(  \theta+\alpha\right)  \right\}    & =-a\\
\theta_{\min}  & =\pi-\alpha
\end{align*}
Hence
\[
F\left(  \mathbf{\varphi,}\left(  \frac{1}{\sqrt{2}},\frac{1}{\sqrt{2}%
}\right)  \right)  =\left\langle 1\left\vert h_{1}1\right.  \right\rangle
+\left\langle 2\left\vert h_{1}2\right.  \right\rangle +\left(  \left.
11\right\vert 11\right)  +\left(  \left.  22\right\vert 22\right)
-2\left\vert \left(  \left.  11\right\vert 22\right)  \right\vert
\]



\end{document}